\section{RL 次单调性与近端点收敛}
\label{sec:rleb-ppa}

沿用第一章的有限维空间 $X=\R^n$、原关系 $F:X\rightrightarrows X$
和步长 $\lambda>0$。固定非空闭目标 $S\subset\zer F$，记
\[
 d(x)=\dist(x,S),\qquad
 P_S(x)=\{p\in S:\|x-p\|=d(x)\},\qquad
 r_F(y)=\dist(0,F(y)).
\]
本章把近端一步的内积估计与真实残差误差界结合起来，
先证明普通 PPA 的局部有限长度，再讨论幂次模、线性能量及固定核推广。

\subsection{普通 PPA 的同窗条件}

\begin{assumption}[普通 PPA 的 RLEB 条件]
\label{ass:rleb-ppa}
取开集 $U\subset X$、$R,\bar t>0$ 和 $0<\kappa<1$，令
\[
 U_R=\{x\in U:d(x)\le R\}.
\]
取有限非减函数 $\omega:[0,R]\to[0,\infty)$ 和
$\psi:[0,\bar t]\to[0,\infty)$，二者在零处连续且
$\omega(0)=\psi(0)=0$。定义
\begin{equation}
\label{eq:ppa-length-modulus}
 B(t)=\frac{t+\omega(t)}2,\qquad
 \Theta(t)=\sum_{j=0}^{\infty}B(\kappa^jt),\quad 0\le t\le R.
\end{equation}
同一个允许近端关系 $\mathcal J$ 满足以下条件。
\begin{enumerate}
\item[(P1)] 对每个 $x\in U_R$，有
$\varnothing\ne\mathcal J(x)\subset J_{\lambda F}(x)$。
\item[(P2)] 对每个 $x\in U_R$ 和每个 $y\in\mathcal J(x)$，
存在 $p\in P_S(x)$，使
\begin{equation}
\label{eq:ppa-zero-anchor}
 \langle y-p,x-y\rangle
 \ge\frac{d(x)^2-\omega(d(x))^2}{4}.
\end{equation}
\item[(P3)] 对每个 $y\in\mathcal J(U_R)$，若 $r_F(y)\le\bar t$，则
\begin{equation}
\label{eq:ppa-true-output-EB}
 d(y)\le\psi(r_F(y)).
\end{equation}
此外，输出评价窗满足
\begin{equation}
\label{eq:ppa-output-window}
 \frac{B(R)}\lambda\le\bar t.
\end{equation}
\item[(P4)] 对全部 $0<t\le R$，有
\begin{equation}
\label{eq:direct-compatibility}
 \psi(B(t)/\lambda)\le\kappa t,
 \qquad \Theta(R)<\infty.
\end{equation}
\end{enumerate}
\end{assumption}

(P1) 是输入覆盖条件；它既不由原关系处处有定义推出，也不由图块近端的单值性推出。
(P2) 中的最近零点可以依赖于实际输出，条件须逐一覆盖 $\mathcal J(x)$ 的全部元素。
若同一图块满足全对 RL，且每个输入有属于该图块的最近零锚，
将输出图点与该零锚配对即可得到 (P2)。这一验证只涉及该图块的输出；
要控制完整 PPA，应取 $\mathcal J=J_{\lambda F}$ 并核对全部完整输出。

(P3) 是定义~\ref{def:true-EB} 在实际输出集上的误差界，评价的是完整 $F$ 的真残差。
它与窗条件~\eqref{eq:ppa-output-window} 分别陈述，以便先验证评价域，再使用误差界。
(P4) 的两项条件分别负责距离收缩和位移可和。
以下将这些零点锚几何、真实残差误差界及兼容可和条件合称为 RLEB 条件。

\subsection{普通 PPA 的局部有限长度}

\begin{theorem}[全部允许选择的局部收敛]
\label{thm:selection-RLEB}
\label{thm:rleb-ppa-general}
在假设~\ref{ass:rleb-ppa} 下，取 $x^0\in U_R$，记 $d_0=d(x^0)$。
若
\begin{equation}
\label{eq:direct-localization}
 \Theta(d_0)<\dist(x^0,X\setminus U),
\end{equation}
则按 $x^{k+1}\in\mathcal J(x^k)$ 逐步选取的每条轨道均可无限延拓；
其全部点属于 $U_R$，轨道具有有限长度，并收敛到某个 $x^\infty\in S\cap U$。
对所有 $k\ge0$，有
\begin{align}
 d(x^k)&\le\kappa^kd_0,\label{eq:ppa-distance-rate}\\
 \|x^{k+1}-x^k\|&\le B(d(x^k))\le B(\kappa^kd_0),\label{eq:ppa-step-rate}\\
 \sum_{k=0}^{\infty}\|x^{k+1}-x^k\|&\le\Theta(d_0),\label{eq:ppa-total-length}\\
 \|x^\infty-x^k\|&\le\Theta(d(x^k))\le\Theta(\kappa^kd_0).
 \label{eq:ppa-tail-bound}
\end{align}
这些估计统一于全部允许选择，极限点可以随选择改变。
特别地，当 $\mathcal J=J_{\lambda F}$ 时，结论覆盖完整近端关系的全部输出。
\end{theorem}

\begin{proof}
任取 $x\in U_R$ 及 $y\in\mathcal J(x)$，取 (P2) 中的最近零点 $p$，
记 $d=d(x)$、$s=\|x-y\|$。引理~\ref{lem:one-step} 给出
\begin{equation}
\label{eq:ppa-one-step}
 d(y)^2+s^2\le\frac{d^2+\omega(d)^2}{2},\qquad s\le B(d).
\end{equation}
由真实近端包含，$(x-y)/\lambda\in F(y)$，从而
\[
 r_F(y)\le\frac{s}\lambda\le\frac{B(d)}\lambda
 \le\frac{B(R)}\lambda\le\bar t.
\]
因此 (P3) 可以用于这个输出。当 $d>0$ 时，(P4) 及 $\psi$ 非减给出
\[
 d(y)\le\psi(r_F(y))\le\psi(B(d)/\lambda)\le\kappa d.
\]
当 $d=0$ 时，式~\eqref{eq:ppa-one-step} 强制 $s=0$，故 $y=x$，
同一距离递推仍然成立。

(P1) 保证初值处存在允许输出。假设已构造到 $x^k\in U_R$，
并已对该有限前缀建立距离递推。任取下一允许输出 $x^{k+1}$，上述一步估计给
\[
 d(x^{k+1})\le\kappa^{k+1}d_0,\qquad
 \|x^{k+1}-x^0\|
 \le\sum_{j=0}^{k}B(\kappa^jd_0)
 \le\Theta(d_0)<\dist(x^0,X\setminus U).
\]
故 $x^{k+1}\in U$，且 $d(x^{k+1})\le R$。于是 (P1) 可以继续使用。
归纳证明了每个有限前缀均留在 $U_R$，并且任一选择都可无限延拓。

现在距离和位移满足定理~\ref{thm:distance-length} 的条件，
取其中的位移包络为 $B$，便得有限长度、点收敛及尾界。
也可直接由步长可和得到 Cauchy 性；$S$ 闭且 $d(x^k)\to0$，
故极限属于 $S$。最后，
\[
 \|x^\infty-x^0\|\le\Theta(d_0)<\dist(x^0,X\setminus U)
\]
使极限仍属于 $U$。所有估计均对每个允许输出成立，因而常数与选择无关。
\end{proof}

\begin{remark}[相对状态集]
\label{rem:relative-state}
上述定理也适用于不变状态集 $D\subset X$：
将假设~\ref{ass:rleb-ppa} 的输入域换成 $D\cap U_R$，
要求全部允许输出仍属于 $D$，初值取自 $D\cap U_R$，
并取非空闭目标 $S=\zer F\cap D$。
有限前缀证明不变，故得到上述全部结论，且 $x^k,x^\infty\in D$。
这里不要求 $D$ 本身闭；极限属于 $D$ 是因为它属于闭目标 $S\subset D$。
\end{remark}

\subsection{幂次模的显式条件与尾界}

\begin{corollary}[幂次 RLEB 条件]
\label{cor:power-RLEB}
在假设~\ref{ass:rleb-ppa} 中取
$\omega(t)=Lt^\gamma$、$\psi(s)=cs^q$，其中
$L\ge0$、$0<\gamma\le1$、$c,q>0$。
若 (P1)--(P3) 成立，且对全部 $0<t\le R$ 有
\begin{equation}
\label{eq:power-compatibility}
 \frac{c}{(2\lambda)^q}(t+Lt^\gamma)^q\le\kappa t,
\end{equation}
则满足严格初值预算~\eqref{eq:direct-localization} 的全部允许轨道
具有定理~\ref{thm:selection-RLEB} 的结论，其中
\begin{equation}
\label{eq:power-tail}
 B(t)=\frac{t+Lt^\gamma}{2},\qquad
 \Theta(t)=\frac12\left(
 \frac{t}{1-\kappa}+\frac{Lt^\gamma}{1-\kappa^\gamma}\right).
\end{equation}
集合距离的 Q 线性因子至多为 $\kappa$，点误差的 R 线性因子至多为 $\kappa^\gamma$。
\end{corollary}

\begin{proof}
对 $B(\kappa^jt)$ 的两项分别求几何级数，即得
式~\eqref{eq:power-tail}，特别地 $\Theta(R)<\infty$。
式~\eqref{eq:power-compatibility} 正是直接兼容条件，故主定理适用。
距离递推给 Q 线性界；尾界中的 $\kappa^k$ 和 $\kappa^{k\gamma}$
两项给点误差的 R 线性界。
\end{proof}

当 $0<\gamma<1$ 且 $L>0$ 时，直接兼容商为
\begin{equation}
\label{eq:power-critical-ratio}
 \frac{\psi(B(t)/\lambda)}t
 =\frac{c}{(2\lambda)^q}
 t^{\gamma q-1}(t^{1-\gamma}+L)^q.
\end{equation}
因此，$\gamma q>1$ 时该商趋于零，可以缩小半径得到统一收缩；
$\gamma q=1$ 时，小窗收缩成立当且仅当
$cL^q/(2\lambda)^q<1$；$\gamma q<1$ 时该商趋于无穷，
这一直接测试不能成立。最后一种情形本身不否定实际轨道收敛。
$\gamma=1$ 或 $L=0$ 时按式~\eqref{eq:power-compatibility} 直接检验。
由于 $B(R)\to0$，缩小半径还可同时保证输出评价窗。

\subsection{线性误差界的能量收缩}

在线性阶，把步长先估计再代入误差界会舍去平方能量的一部分信息。
保留式~\eqref{eq:ppa-one-step} 的两项可以得到另一充分收缩条件。

\begin{corollary}[线性能量收缩]
\label{cor:linear-energy-convergence}
沿用假设~\ref{ass:rleb-ppa} 的基本对象，仅要求 (P1)--(P3)，
其中 $\omega(t)=Lt$、$\psi(s)=Ks$，$L\ge0$、$K>0$。
定义
\begin{equation}
\label{eq:linear-energy-factor}
 \rho_E=\left(\frac{1+L^2}{2\{1+(\lambda/K)^2\}}\right)^{1/2},
 \qquad \Theta_E(t)=\frac{(1+L)t}{2(1-\rho_E)}.
\end{equation}
若 $\rho_E<1$，且 $x^0\in U_R$ 满足
\begin{equation}
\label{eq:linear-energy-budget}
 \Theta_E(d_0)<\dist(x^0,X\setminus U),
\end{equation}
则主定理的全部结论成立，其中距离因子和尾预算分别换成
$\rho_E$ 和 $\Theta_E$。特别地，
\[
 d(x^k)\le\rho_E^kd_0,\qquad
 \|x^\infty-x^k\|\le\Theta_E(d(x^k))\le\Theta_E(\rho_E^kd_0).
\]
此结论也适用于注~\ref{rem:relative-state} 的相对状态集。
\end{corollary}

\begin{proof}
任取合法转移 $x\mapsto y$，记 $d=d(x)$、$d^+=d(y)$、$s=\|x-y\|$。
引理~\ref{lem:one-step} 和 (P3) 的输出窗先给
\[
 s\le\frac{1+L}{2}d,\qquad
 r_F(y)\le\frac{s}\lambda\le\bar t.
\]
故真实残差误差界可以使用，并与一步能量合为
\[
 (d^+)^2+s^2\le\frac{1+L^2}{2}d^2,
 \qquad d^+\le K r_F(y)\le\frac K\lambda s.
\]
消去 $s^2$ 得
\[
 \{1+(\lambda/K)^2\}(d^+)^2\le\frac{1+L^2}{2}d^2,
\]
即 $d^+\le\rho_Ed$，包括 $d=0$ 的情形。
对已构造的有限前缀，有
\[
 \|x^{k+1}-x^0\|
 \le\frac{1+L}{2}\sum_{j=0}^{k}\rho_E^jd_0
 \le\Theta_E(d_0)<\dist(x^0,X\setminus U).
\]
所以每一步仍在 $U_R$，覆盖条件可再次使用。
无限延拓后求和及从任意指标重启求和，分别给有限长度和所列尾界；
距离趋零、$S$ 闭及严格预算保证极限属于 $S\cap U$。
\end{proof}

为与 Luke--Tam 正式版所列的估计比较，固定同一图块、步长和输出残差窗口，
并要求全部允许输出及相应最近零点锚均属于该图块。
当该图上的 $\lambda F$ 满足 Luke--Tam 的 submonotonicity 条件时，
线性 RL 参数满足 $L^2=1+4\tau$、$\tau\ge0$
\cite[Definition~2]{LukeTam2025GeneralizedMonotonicity}。
于是本推论的能量收缩条件为
\begin{equation}
\label{eq:energy-LT-budget}
 2\tau(K/\lambda)^2<1.
\end{equation}
另一方面，同一真误差界和三角不等式给
\[
 d\le s+d^+\le(1+K/\lambda)s.
\]
在一步能量中代入这个步长下界，即得
\begin{equation}
\label{eq:LT-displayed-factor}
 (d^+)^2\le\kappa_{\rm LT}(K)^2d^2,
 \qquad
 \kappa_{\rm LT}(K)^2
 =1+2\tau-\left(\frac{\lambda}{\lambda+K}\right)^2.
\end{equation}
相应的显示充分条件为
\begin{equation}
\label{eq:LT-displayed-budget}
 \tau(1+K/\lambda)^2<\frac12.
\end{equation}
这给出了正式版所列的充分参数条件与距离因子
\cite[Assumption~2, Theorem~2]{LukeTam2025GeneralizedMonotonicity}。
原证明使用任意略大于 $K$ 的误差界常数 $\rho$；
这里假定 $K$ 已是同一窗口内有效的误差界常数，
因此 $\kappa_{\rm LT}(K)$ 是原文因子在 $\rho\downarrow K$ 时的下确界。

置 $a=K/\lambda>0$，直接计算得
\begin{equation}
\label{eq:LT-energy-factor-gap}
 \kappa_{\rm LT}(K)^2-\rho_E^2
 =\frac{2\{a+\tau(1+a)^2\}}{(1+a^2)(1+a)^2}>0.
\end{equation}
式~\eqref{eq:LT-displayed-budget} 中的 Luke--Tam 充分条件
蕴含式~\eqref{eq:energy-LT-budget}，而反向不成立。
在两者都给出收缩时，能量因子严格更小；同一步长包络下，
$\Theta_E$ 也相应缩小。这些比较使用相同目标、输入覆盖、
允许输出及真实残差评价窗，改善的是上述显示充分条件和因子，
并不由此断言一个完整局部定理包含另一个。
当 $\tau=0$ 时，
\[
 \rho_E=\frac{K}{\sqrt{K^2+\lambda^2}},
\]
恢复 Leventhal 的单调端点因子
\cite[Theorem~3.1]{Leventhal2009MetricSubregularityPPA}。
当 $L<1$ 时，RL 给出更强的正耗散估计；此时不把负数
$(L^2-1)/4$ 赋给 Luke--Tam 定义中的非负缺陷参数。

\subsection{能量因子与长度估计的锐性见证}

先给出同时达到能量因子、步长和总长度估计的标量族。
给定 $0<\kappa<1$ 和 $\lambda>0$，令
\begin{equation}
\label{eq:scalar-sharp-family}
 F(y)=-\frac{1+\kappa}{\lambda\kappa}y,\qquad
 J_{\lambda F}(x)=-\kappa x,\qquad S=\{0\}.
\end{equation}
取 $U=\R$、任意 $R>0$，并置
\[
 L=1+2\kappa,\qquad K=\frac{\lambda\kappa}{1+\kappa},
 \qquad \bar t\ge\frac{(1+\kappa)R}\lambda.
\]
直接求解近端方程可知，对每个输入，唯一输出均为 $y=-\kappa x$，
因此覆盖全域成立。
对于 $y=-\kappa x$，有
\[
 \langle y,x-y\rangle=-\kappa(1+\kappa)|x|^2
 =\frac{1-L^2}{4}|x|^2,
 \qquad r_F(y)=\frac{|y|}{K}.
\]
所以零锚几何和真误差界均取等，且
\[
 |y|^2+|x-y|^2
 =\{\kappa^2+(1+\kappa)^2\}|x|^2
 =\frac{1+L^2}{2}|x|^2.
\]
于是 $\rho_E=\kappa$，步长恰为 $B(|x|)=(1+\kappa)|x|$。
轨道为 $x^k=(-\kappa)^kx^0$，其总长满足
\begin{equation}
\label{eq:scalar-sharp-length}
 \sum_{k=0}^{\infty}|x^{k+1}-x^k|
 =\frac{1+\kappa}{1-\kappa}|x^0|=\Theta_E(|x^0|).
\end{equation}
这说明显示估计在这一参数族上不能一致减小。
这里使用的是算子近端关系；相应凹二次函数并不给出同一个 argmin 近端最小化问题。

平面旋转族则表明能量因子可由非标量完整关系达到。
令 $Q=\begin{pmatrix}0&-1\\1&0\end{pmatrix}$，取
$\alpha,\beta\in\R$，使
$m^2=\alpha^2+\beta^2>0$、$D=1+2\alpha+m^2>0$，并定义
\begin{equation}
\label{eq:rotation-sharp-family}
 F=\frac{\alpha\Id+\beta Q}{\lambda},\qquad
 T=J_{\lambda F}=\frac{(1+\alpha)\Id-\beta Q}{D}.
\end{equation}
有 $r_F(y)=m\|y\|/\lambda$，故真误差界的锐常数为 $K=\lambda/m$。
对任意输入差 $z$，直接计算得到
\[
 \|Tz\|^2=\frac{\|z\|^2}{D},\qquad
 \|z-Tz\|^2=\frac{m^2}{D}\|z\|^2,
 \qquad
 \langle Tz,z-Tz\rangle=\frac{\alpha}{D}\|z\|^2.
\]
因此，当参数取为下式中的 $L$ 时，全对线性 RL 不等式取等，且能量因子为
\begin{equation}
\label{eq:rotation-sharp-parameters}
 L^2=\frac{1-2\alpha+m^2}{1+2\alpha+m^2},
 \qquad \rho_E^2=\frac1D.
\end{equation}
当 $D>1$ 时，这一族的完整近端迭代从任意初值出发均收敛到零点，
并达到上述能量距离因子。
若预先指定 $m=\lambda/K>0$ 与 $L\ge0$，则它的可行范围恰为
\begin{equation}
\label{eq:rotation-feasible-range}
 \left|\frac{(1-L^2)(1+m^2)}{2(1+L^2)}\right|\le m;
\end{equation}
此时取括号内的数为 $\alpha$，再取
$\beta^2=m^2-\alpha^2$ 即可。这个范围只陈述旋转族的达到性。

例如取 $\lambda=K=1$、$\alpha=-2/7$、$\beta=-3\sqrt5/7$，则
\[
 F=-\frac27\Id-\frac{3\sqrt5}{7}Q,\qquad
 J_F=\frac12\Id+\frac{3\sqrt5}{10}Q,
 \qquad L^2=\frac95,\quad \tau=\frac15,\quad\rho_E^2=\frac7{10}.
\]
令 $U=\R^2$、$R=1$、$\bar t=2$，全部输出的真残差误差界与评价窗均成立。
此时，Luke--Tam 充分条件~\eqref{eq:LT-displayed-budget}
的左端为 $4/5>1/2$，故该条件不成立；
能量收缩条件~\eqref{eq:energy-LT-budget}
的左端为 $2/5<1$，故该条件成立。
这给出同一完整关系上的严格参数差异。
能量收缩条件在所有参数范围内并非必要条件；相同前提还给出直接因子
$K(1+L)/(2\lambda)$，而旋转达到性也受式~\eqref{eq:rotation-feasible-range} 限制。

\subsection{完整多值选择与非 Lipschitz 漂移}

\paragraph{完整多值近端关系.}
\label{ex:complete-multivalued}
在 $\R$ 上令
\[
 F(y)=\{2y,3y\},\qquad \lambda=1,\qquad S=\{0\}.
\]
两条近端方程穷尽完整输出，给出
\[
 J_F(x)=\{x/3,x/4\}.
\]
取 $U=\R$、$R=\bar t=1$、$\omega(t)=t/2$、
$\psi(s)=s/2$ 和 $\kappa=3/8$。
最近零点为 $p=0$；若 $y=x/3$，则 $y(x-y)=2x^2/9$，
若 $y=x/4$，则 $y(x-y)=3x^2/16$。
两者均不小于 $\{x^2-(|x|/2)^2\}/4=3x^2/16$，故 (P2) 对全部输出成立。
完整真残差为 $r_F(y)=2|y|$，所以真误差界取等。
此外，
\[
 B(R)=\frac34\le\bar t,\qquad
 \psi(B(t))=\frac38t,
\]
可和性由推论~\ref{cor:power-RLEB} 给出，留域预算自动成立。
因此任意切换均由定理~\ref{thm:selection-RLEB} 控制。
同一非零输入的两个输出不同，而完整全对 RL 在零输入差时必须使输出相同；
故这个完整关系不满足任何零处消失模的全对 RL。
例子说明逐输出零锚条件确实允许严格多值更新。

\paragraph{法向收缩与切向平方根漂移.}
\label{ex:sqrt-tangential-drift}
在 $\R^2$ 上取完整单值原关系
\begin{equation}
\label{eq:sqrt-triangular-operator}
 F(u,v)=(u,-\sqrt{|u|}),\qquad
 S=\{0\}\times\R,\qquad \lambda=1.
\end{equation}
近端方程的第一坐标给 $u=a/2$，第二坐标再唯一确定 $v$，因而
\begin{equation}
\label{eq:sqrt-triangular-resolvent}
 J_F(a,b)=(a/2,b+\sqrt{|a|/2}).
\end{equation}
完整真残差满足
\[
 r_F(u,v)^2=u^2+|u|,\qquad
 d((u,v),S)=|u|\le r_F(u,v)^2.
\]
取 $U=\R^2$、$R=1/16$、$\bar t=1$、
$\omega(t)=\sqrt{2t}$、$\psi(t)=t^2$、$\kappa=3/4$。
输入 $(a,b)$ 的最近零点唯一，为 $p=(0,b)$。
对式~\eqref{eq:sqrt-triangular-resolvent} 的实际输出 $y$，有
\[
 \langle y-p,(a,b)-y\rangle
 =\frac{a^2}{4}-\frac{|a|}{2}
 =\frac{d((a,b),S)^2-\omega(d((a,b),S))^2}{4}.
\]
故锚内积恰好达到下界。评价窗和直接兼容分别为
\[
 B(R)=\frac{1+4\sqrt2}{32}<1,\qquad
 \sup_{0<t\le R}\frac{\psi(B(t))}{t}
 =\frac{33+8\sqrt2}{64}<\frac34.
\]
由推论~\ref{cor:power-RLEB}，这些条件控制完整轨道。
实际迭代可以进一步求和为
\begin{equation}
\label{eq:sqrt-triangular-orbit}
 a_k=2^{-k}a_0,\qquad
 b_k=b_0+\frac{\sqrt{|a_0|}}{\sqrt2-1}(1-2^{-k/2}),
\end{equation}
所以
\[
 x^\infty=\left(0,b_0+\frac{\sqrt{|a_0|}}{\sqrt2-1}\right).
\]
平方根差估计
$|\sqrt{|a|}-\sqrt{|a'|}|\le\sqrt{|a-a'|}$
说明近端映射局部为 $1/2$ 阶 H\"older 连续。
比较输入 $(a,b)$ 与 $(0,b)$ 可知，任何高于 $1/2$ 的局部指数均不成立。
同样，$r_F(u,v)^2\sim|u|$（$u\to0$）说明
$d\le c r_F^q$ 的最大可用幂次为 $q=2$。
当 $a_0\ne0$ 时，集合距离的实际因子为 $1/2$，而
\[
 \|x^k-x^\infty\|^2
 =4^{-k}a_0^2+
 \frac{|a_0|}{(\sqrt2-1)^2}\,2^{-k},
\]
故点误差的渐近比值为 $1/\sqrt2$。
法向距离与切向位移处于不同的幂次尺度，因而距离收缩率与点收敛率可以不同。
